\documentclass[11pt]{article}
\usepackage[margin=1in]{geometry}
\usepackage{amsmath,amssymb,amsthm,booktabs,array}
\usepackage[T1]{fontenc}
\usepackage{lmodern}
\usepackage{xeCJK}
\setCJKmainfont{FandolSong-Regular.otf}
\usepackage{microtype}
\usepackage[hidelinks]{hyperref}
\setlength{\emergencystretch}{2em}
\newtheorem{theorem}{Theorem}
\newtheorem{lemma}{Lemma}
\newcommand{\R}{\mathbb R}
\newcommand{\ext}{\operatorname{ext}}
\newcommand{\code}[1]{\texttt{\detokenize{#1}}}
\title{Extreme points are not dense in a tensor nuclear-norm ball\\
\large A counterexample to the ordinary density clause of conjecture 00000007302}
\author{C0ldSmi1e}
\date{October 12, 2026}
\begin{document}
\maketitle
\begin{abstract}
In the real tensor product of three Euclidean two-dimensional spaces, equip tensors
with the nuclear norm defined by finite rank-one decomposition costs. The origin
belongs to its closed unit ball but lies outside the closure of the ball's extreme
points. Thus those extreme points are not dense relative to the ball. The Lean
proof constructs the actual tensor nuclear norm, proves its positive definiteness,
and establishes the relative non-density conclusion. This refutes the ordinary
density assertion in the supplied conjecture; its additional, undefined
parameter-space clause is not reconstructed or independently formalized.
\end{abstract}

\section{Original statement and interpretation}
The exact English statement is:
\begin{quote}
Definition: The tensor nuclear norm as the convex envelope of low-rank tensors.
Conjecture: The density of extreme points of the nuclear-norm ball holds, and the
combination of the density is the parameter space of rank-one decompositions.
(rank-one-parameter nuclear-ball extreme density)
\end{quote}
The exact Chinese statement is:
\begin{quote}
定义：张量核范数：低秩张量的凸包络。猜想：核范数的球：核范数球的极点的稠密且密的组合为秩一分解的参数空间。（张量核范数秩一参数空间分解极点稠密球）
\end{quote}
The unchanged bilingual source is also supplied as \code{conjecture.md}.

The text does not specify the tensor dimensions, scalar field, topology, or
the meaning of its parameter-space identification. The phrase about a convex
envelope alone is not a complete norm definition. We therefore state our
conventions explicitly: real Euclidean factors, the usual finite-decomposition
nuclear (projective) norm, the centered closed unit ball, real convex extreme
points, and density relative to that ball in its norm topology.

The English sentence separately asserts density of the extreme points and joins
the remaining clause by ``and.'' Under this ordinary reading, the false density
assertion already refutes the conjunction. Neither language explicitly replaces
the extreme points by their convex combinations or the ball by its boundary.
We do not claim to settle those different assertions, or to provide a unique
formal translation of the entire underspecified sentence. A later clarification
of a different intended density target would require reconsidering the scope.

\section{The genuine tensor nuclear norm}
Let $H=\R^2$ with its Euclidean norm and let
\[
 T=H\otimes_{\R}H\otimes_{\R}H.
\]
More precisely, the implementation uses the algebraic tensor product of the
family indexed by $\{0,1,2\}$. For $x\in T$, define
\begin{equation}\label{eq:nuclear}
 \nu(x)=\inf\left\{
 \sum_{j=1}^{N}\|u_j\|_2\|v_j\|_2\|w_j\|_2:
 x=\sum_{j=1}^{N}u_j\otimes v_j\otimes w_j,\quad N\ge0
 \right\}.
\end{equation}
The infimum ranges over arbitrary finite decompositions, with no rank bound or
restriction to the coordinate decomposition. Empty lists are allowed. Every
tensor has a finite decomposition, and every cost is nonnegative; hence the
feasible cost set is nonempty and bounded below.

Mathlib's \code{PiTensorProduct.projectiveSeminorm} uses finite weighted sums
$\sum_j c_j\,u_j\otimes v_j\otimes w_j$, with cost
$\sum_j |c_j|\|u_j\|_2\|v_j\|_2\|w_j\|_2$. The formal bridge sends such a
term to $(c_ju_j)\otimes v_j\otimes w_j$. This preserves both the tensor and its
cost. Conversely, an ordinary decomposition is a weighted one with all
coefficients equal to one. The two feasible cost sets are therefore equal,
so the library seminorm equals \eqref{eq:nuclear}.

The library supplies nonnegativity, the triangle inequality, and scalar
homogeneity. The remaining point, proved in this submission, is positive
definiteness on the actual tensor product.

\begin{lemma}\label{lem:definite}
If $\nu(x)=0$ for $x\in T$, then $x=0$.
\end{lemma}
\begin{proof}
Let $\phi:T\to\R$ be any algebraic linear functional and put
$f(u,v,w)=\phi(u\otimes v\otimes w)$. Every trilinear form on $H^3$ is
continuous: writing $e_0,e_1$ for the Euclidean basis, its explicit expansion is
\[
 f(m_0,m_1,m_2)=
 \sum_{a\in\{0,1\}^{\{0,1,2\}}}
 \left(\prod_{i=0}^{2}(m_i)_{a_i}\right)
 f(e_{a_0},e_{a_1},e_{a_2}).
\]
This is a finite sum of continuous coordinate products. For every decomposition
in \eqref{eq:nuclear}, the multilinear operator-norm bound gives
\[
 |\phi(x)|\le \|f\|\sum_j\|u_j\|_2\|v_j\|_2\|w_j\|_2.
\]
Taking the infimum yields $|\phi(x)|\le\nu(x)\|f\|=0$. Thus every algebraic
linear functional vanishes at $x$. Algebraic dual separation implies $x=0$.
Both the continuity expansion and this argument are formalized.
\end{proof}

The resulting definite seminorm is installed as the normed additive group
structure on $T$, and its scalar homogeneity gives a real normed-space
structure. The equality between the installed norm and $\nu$ is a proved
theorem, not an assumption. Every closure below uses precisely this norm
topology; no unrelated array or ambient norm is substituted.

The tensor $q=e_0\otimes e_0\otimes e_0$ is nonzero: the linear lift of the
trilinear form $(u,v,w)\mapsto u_0v_0w_0$ takes $q$ to $1$ and the zero tensor to $0$.
This supplies the nontriviality required in the geometric argument.

\section{The obstruction to density}
For a convex set $B$ in a real vector space, an extreme point $x\in B$ is a
point such that if $x$ lies in the open segment between $a,b\in B$, then
$a=b=x$. We use Mathlib's actual \code{Set.extremePoints} predicate.

\begin{lemma}\label{lem:sphere}
In any nontrivial real normed space, every extreme point of the closed unit
ball belongs to the unit sphere.
\end{lemma}
\begin{proof}
Suppose $\|x\|<1$. Choose a unit vector $e$ and let
$y=(1-\|x\|)e/2$. Then $y\ne0$ and
\[
 \|x+y\|,\ \|x-y\|
 \le\|x\|+\|y\|=\frac{1+\|x\|}{2}<1.
\]
The point $x$ is the midpoint of the two distinct ball points $x+y$ and $x-y$,
so it is not extreme. Every extreme point already has norm at most one,
and therefore has norm exactly one.
\end{proof}

\begin{theorem}\label{thm:counterexample}
For the nuclear ball $B=\{x\in T:\nu(x)\le1\}$,
\[
 0\in B,\qquad 0\notin\overline{\ext(B)},\qquad
 B\not\subseteq\overline{\ext(B)}.
\]
In particular, the extreme points are not dense relative to $B$.
\end{theorem}
\begin{proof}
Lemma~\ref{lem:sphere} puts all extreme points on the unit sphere of the
nuclear norm. That sphere is closed, so it contains the closure of those
extreme points. The zero tensor belongs to the ball and has norm zero,
so it is outside that closure.
\end{proof}

For any $S\subseteq B$, density in the relative topology on $B$ means
$B\subseteq\overline S$, where the closure on the right is in $T$.
Thus the displayed negation addresses relative density in the ball, not just
density in the whole tensor space. It does not require classifying the
extreme points or proving uniqueness of rank-one decompositions.

\section{Formal statements and verification}
The project is pinned to Lean 4.19.0 and Mathlib v4.19.0, revision
\begin{center}\small\code{c44e0c8ee63ca166450922a373c7409c5d26b00b}.\end{center}
The files have the following roles:
\begin{itemize}
\item \code{TensorNuclear.lean}: actual Euclidean factors and tensor product,
continuity of trilinear forms, positive definiteness, installed nuclear norm,
and a concrete nonzero tensor.
\item \code{DecompositionBridge.lean}: cost-preserving conversion between
weighted and ordinary finite decompositions, existence and lower bounds,
and equality of their infima.
\item \code{GenericExtreme.lean}: Lemma~\ref{lem:sphere} and its general
relative non-density consequence.
\item \code{Conjecture7302.lean}: the explicit nuclear norm and ball,
their equality with the installed norm and metric ball, and
Theorem~\ref{thm:counterexample}.
\end{itemize}
\clearpage
The final declarations in namespace \code{Conjecture7302} are
\begin{center}\small
\code{nuclear_extremePoints_not_relatively_dense}\\[3pt]
\code{nuclear_ball_counterexample}.
\end{center}
The former proves
\[
 \neg\bigl(B\subseteq\operatorname{closure}(B.\mathrm{extremePoints}\ \R)\bigr),
\]
and the latter supplies the origin as the explicit witness.

From the supplied \code{lean} directory, the ordinary build command is
\code{lake build}. The auxiliary \code{AxiomAudit.lean} prints the dependency
axioms of the principal declarations and can be replayed with
\code{lake env lean -DwarningAsError=true AxiomAudit.lean}.
The accompanying \code{VERIFICATION.md} records the separate author and
independent checks, exact source hashes, complete source replay, and dependency
audit. This analytic proof requires no numerical search or auxiliary numerical
calculation. No admitted proofs, custom axioms, or \code{native_decide} are used.
The permitted foundational axioms are only \code{propext},
\code{Classical.choice}, and \code{Quot.sound}.

Dependency analysis and kernel checking have different roles: the audit
inspects declaration dependencies, while source compilation checks proofs
with Lean's kernel. Cached library imports are pinned and checked; the whole
Mathlib library and the Lean runtime are not rebuilt from bootstrap. Local
validation is distinct from maintainer acceptance. These engineering checks
also do not remove the interpretation limits stated in Section~1.
\end{document}
